\section{Experimental Evaluation}

Every figure in this section comes from one host in one pinned sitting on 6~September 2026: an AMD Ryzen~7 5800H, 16 logical CPUs, Linux 7.2.3, clock pinned to 3.2\,GHz with boost disabled and the setting verified by read-back against the state captured before pinning. The timed work ran in the project's \texttt{ctpb-dev:2026-09-06} image built from the repository's own \texttt{Dockerfile} \cite{ctpb2026}, compiled with Clang 21.0 (P2996 branch) under \texttt{-std=c++26 -freflection -O3} against Protobuf 3.21.12. Every raw trial line, the clock samples taken during the run and the exact benchmark script are published with the paper.
% src: measurements/2026-09-06-laptop-p5-single-host/README.md conditions table (host, image,
%   toolchain, clock: 41 samples across 16 CPUs, median 3200 MHz, min 3159, spread 0.97%);
%   raw/host-state.txt (BEFORE PIN / DURING / AFTER RESTORE, verification=ok against the
%   captured baseline rather than a hardcoded value)

The numbers below replace those of the previous version rather than extending them, and the reason is worth stating. The earlier table was taken on a different machine, an AMD Ryzen~5 3600, under a harness this repository has since repudiated for three named defects, any one of which could account for the whole reported margin. That harness also decoded into a freshly constructed destination on the \texttt{ctpb} arm while reusing one on the reference arm, charging \texttt{ctpb} a heap allocation per iteration that its competitor did not pay. Re-measured on one host with that asymmetry removed, the encode margin roughly doubles and the decode margin falls well below what was previously printed. No \texttt{ctpb} code changed between the two campaigns; the instrument did.
% src: rerun_all.sh (the repudiation and the three defects); measurements/2026-09-06-laptop-p5-single-host/README.md
%   "What moved" and "Two instrument defects found and fixed before these numbers were read";
%   method/diff-from-committed-eval.txt is the exact change from the committed harness

A single sitting on one host is one sitting on one host. These cells carry no validity gates, because the scripts that produce them have none, and nothing here was accepted or refused by a rule. What the campaign buys is that every ratio comes from one binary per arm measured minutes apart under a clock that was held and checked, which the table it replaces did not have.
% src: README.md conditions table, gates row ("none -- these scripts have no validity gates,
%   so nothing was accepted or refused")

\subsection{Compile-Time Schema Parsing}
We measured the \texttt{constexpr} parser on three schemas scaled by field count alone (5, 25, and 100 \texttt{int32} fields), each over $20{,}000$ iterations after a $2{,}000$-iteration warmup, across seven trials (Table~\ref{tab:compile_time_results}). Trial~0 is excluded as warm-up throughout this section and the exclusion is declared because it is not cosmetic: on the runtime arms the reference implementation runs 4 to 6\,\% slow on trial~0 while \texttt{ctpb} moves by under $0.4\,\%$, so counting it would flatter \texttt{ctpb}. Spreads are the full range of the six retained trials rather than a standard deviation, because six points do not support one.

\begin{table}[ht]
\centering
\caption{Compile-Time Parser Latency, median and range of six trials}
\label{tab:compile_time_results}
\renewcommand{\arraystretch}{1.15}
\setlength{\tabcolsep}{3pt}
\scriptsize
\begin{tabular}{|l|c|c|c|}
\hline
\textbf{Schema Size} & \textbf{Fields} & \textbf{Latency / Parse} & \textbf{Throughput} \\ \hline
Small & 5 & $0.646\,\mu\text{s}$ [$0.643$, $0.648$] & $1{,}547{,}988\,\text{parses/s}$ \\ \hline
Medium & 25 & $2.515\,\mu\text{s}$ [$2.511$, $2.519$] & $397{,}614\,\text{parses/s}$ \\ \hline
Large & 100 & $9.465\,\mu\text{s}$ [$9.438$, $9.484$] & $105{,}652\,\text{parses/s}$ \\ \hline
\end{tabular}
\end{table}
% src: measurements/2026-09-06-laptop-p5-single-host/raw/compile_time_latency.txt, trials 1-6,
%   median and min/max per column. Throughput is the reciprocal of the median, derived here
%   rather than recorded.

A 100-field schema is ingested in $9.465\,\mu\text{s}$, comfortably under the $11.4\,\mu\text{s}$ this work has claimed for it. The end-to-end cost of invoking \texttt{protoc} was measured only under the repudiated harness and was not re-taken in this campaign, so the ratio between the two is not restated here.
% src: README.md compile-time table ("Under the paper's below-11.4 us claim"). The 7 ms protoc
%   figure and the ~620x-8,150x reduction derived from it appear in no record of this campaign;
%   combining a repudiated figure with a current one would put two instruments in one ratio.

The parser's ceiling is a property of the compiler rather than of the schema language. Bisection at Clang's default \texttt{constexpr} step budget puts the largest single message that compiles at 300 fields, with 301 failing; raising the budget to $4{,}194{,}304$ steps compiles both 500 and 1000 fields. Compiling with the documented default stated explicitly fails in the same place, which is consistent with that default being $1{,}048{,}576$, and the project sets no \texttt{constexpr} step flag anywhere in its own build.
% src: raw/constexpr_step_limit.txt sections 1-4: bisection (375, 312 STEP_LIMIT; 281, 296, 300
%   pass; 304, 302, 301 STEP_LIMIT); -fconstexpr-steps=4194304 rescues 500 and 1000; explicit
%   default also fails; no constexpr flag anywhere in the build

\subsection{Runtime Serialization Throughput}
\label{sec:runtime_comparison}
We evaluated encoding and decoding on a representative 3-field \texttt{Person} message (47-byte wire payload) over seven trials of $10^6$ iterations each with a $20{,}000$-iteration warmup, against both the full \texttt{libprotobuf} and \texttt{libprotobuf-lite} (\texttt{optimize\_for = LITE\_RUNTIME}). Both references were built and run in the same sitting from the same source, so the two columns of Table~\ref{tab:runtime_comparison} are comparable with each other as well as with \texttt{ctpb}.

\begin{table}[ht]
\centering
\caption{Runtime latency, median and range of six trials, 47-byte payload}
\label{tab:runtime_comparison}
\renewcommand{\arraystretch}{1.15}
\setlength{\tabcolsep}{2pt}
\tiny
\begin{tabular}{|l|c|c|c|}
\hline
\textbf{Operation} & \textbf{Latency (ns)} & \textbf{Throughput} & \textbf{Speedup} \\ \hline
\textbf{\texttt{ctpb} encode} & $\mathbf{8.27}$ [$8.24$, $8.51$] & $\mathbf{120.9\,\text{MOps/s}}$ & baseline \\ \hline
\texttt{libprotobuf-lite} encode & $46.11$ [$45.97$, $46.87$] & $21.7\,\text{MOps/s}$ & $\mathbf{5.572\times}$ \\ \hline
\texttt{libprotobuf} encode & $45.64$ [$45.62$, $45.65$] & $21.9\,\text{MOps/s}$ & $\mathbf{5.334\times}$ \\ \hline\hline
\textbf{\texttt{ctpb} decode} & $\mathbf{76.52}$ [$76.08$, $77.73$] & $\mathbf{13.1\,\text{MOps/s}}$ & baseline \\ \hline
\texttt{libprotobuf-lite} decode & $79.83$ [$79.56$, $81.04$] & $12.5\,\text{MOps/s}$ & $\mathbf{1.043\times}$ \\ \hline
\texttt{libprotobuf} decode & $87.80$ [$87.78$, $88.07$] & $11.4\,\text{MOps/s}$ & $\mathbf{1.150\times}$ \\ \hline
\end{tabular}
\end{table}
% src: measurements/2026-09-06-laptop-p5-single-host/raw/runtime_prealloc.txt, both blocks,
%   trials 1-6. The ctpb decode cell is the lite-arm run; the full-libprotobuf arm's own ctpb
%   decode median is 76.33 and its encode 8.56, which is why the two speedups are quoted
%   against their own arm's ctpb median rather than against one shared cell. Speedups are
%   computed from unrounded medians and do not reproduce exactly from the two-decimal cells.
%   Throughput is the reciprocal of the median, derived here rather than recorded.

Encoding is where the compile-time approach pays. Field sizes are known to the compiler, so the encoder allocates one exactly sized buffer through \texttt{byte\_size()} and writes through raw pointers with no reallocation and no virtual dispatch; the margin is $5.57\times$ against the lite runtime and $5.33\times$ against the full one. That the two agree closely is the point of running both: the encode margin does not depend on the descriptor pools, reflection registries and arenas that only the full build carries.

Decoding is a much narrower result and carries a condition. Correcting a per-iteration allocation in the benchmark moves \texttt{ctpb}'s decode from $0.98\times$ to $1.04\times$ against \texttt{libprotobuf-lite}; \texttt{ctpb}'s decoder omits the UTF-8 validation the reference implementation performs, so the remaining margin is not a like-for-like comparison. proto3 requires \texttt{string} fields to be validated as UTF-8, the lite runtime spends about seven per cent of its decode path doing it, and \texttt{ctpb} does not do it at all. The decode figure should be read as a rough parity that still has a correctness asymmetry inside it, pointing the other way. The encode and footprint results carry no such caveat.
% src: README.md "the sentence that has to travel with it" (0.98x -> 1.04x on removing the
%   per-iteration allocation; no ctpb code changed, decode_from fills an existing object and
%   string fields go through assign()); the 7.37 % inclusive figure for IsStructurallyValidUTF8 /
%   VerifyUtf8String / VerifyUTF8 in libprotobuf-lite, and no utf8 anywhere in ctpb's headers

\subsection{Binary Footprint}
\begin{table}[ht]
\centering
\caption{Binary Footprint and Code Generation, exact bytes}
\label{tab:binary_footprint}
\renewcommand{\arraystretch}{1.1}
\setlength{\tabcolsep}{3pt}
\scriptsize
\begin{tabular}{|l|c|c|c|}
\hline
\textbf{Metric} & \textbf{\texttt{ctpb}} & \textbf{\texttt{protoc} (Dyn)} & \textbf{\texttt{protoc} (Static)} \\ \hline
Generated Source Lines & \textbf{0} & 754 & 754 \\ \hline
Stripped Executable & \textbf{18{,}776\,B} & 27{,}504\,B & 1{,}759{,}368\,B \\ \hline
Text Section (\texttt{.text}) & \textbf{10{,}047\,B} & 16{,}373\,B & 1{,}731{,}945\,B \\ \hline
External Shared Libraries & \textbf{0} & \texttt{libprotobuf.so} & \textbf{0} \\ \hline
\end{tabular}
\end{table}
% src: measurements/2026-09-06-laptop-p5-single-host/raw/binary_size.txt sections 1 (379 + 375
%   = 754 lines against 0), 4 ("Exact sizes (bytes)"), 5 (size(1) text column of the stripped
%   executables) and 6 (ldd: ctpb links only libstdc++, libm, libgcc_s and libc).
%   The object-file row of the previous version is dropped rather than restated: the record
%   holds object sizes only as ls -lh rounding (22K, 26K, 5.5K), which is where the earlier
%   21.9 and 31.3 KB tenths came from, and a tenth of a kilobyte is not recoverable from them.

For a single 3-field message \texttt{protoc} emits 754 lines across \texttt{.pb.h} and \texttt{.pb.cc}, whereas \texttt{ctpb} emits no files at all (Table~\ref{tab:binary_footprint}). In fully static deployments \texttt{ctpb}'s stripped executable is $18{,}776$ bytes against $1{,}759{,}368$, a $93.70\times$ reduction, because \texttt{libprotobuf} statically bundles descriptor pools, reflection registries, mutex tables, and arena allocators while \texttt{ctpb} compiles to exact inlined field loops. Against the dynamic build the reduction is $1.46\times$, and the dynamic \texttt{ctpb} build links only \texttt{libc}, \texttt{libm}, \texttt{libstdc++} and \texttt{libgcc\_s}, removing \texttt{libprotobuf.so} version skew from deployment. These are byte counts rather than rounded kilobytes, which is what makes the ratio checkable: the figure this work has claimed all along is right, and it could not previously be recovered from what was published with it.
% src: raw/binary_size.txt section 4; 1759368 / 18776 = 93.6996, and 27504 / 18776 = 1.4649

\subsection{Reference Interoperability and gRPC Parity}
\label{sec:interop}
Correctness was established by dual-direction interoperability testing against \texttt{libprotobuf} 3.21.12 and the \texttt{protoc} code generator \cite{protobuf_spec}. Messages serialized by Google's \texttt{SerializeToString()} decode through \texttt{ctpb::decode<T>()} with full fidelity across scalars, Zigzag integers, strings, enums, nested messages, \texttt{bytes} fields, and packed repeated scalars; in the reverse direction, buffers produced by \texttt{ctpb::encode()} parse cleanly through Google's \texttt{ParseFromArray()}, with all wire tags, varints, and length delimiters conforming to reference parser expectations. Packed repeated fields and \texttt{bytes} are covered explicitly because both are encoded differently from the singular scalar path and are the cases where an independent implementation most easily diverges.

Transport parity was verified against gRPC \cite{grpc_spec}, whose HTTP/2 framing prefixes each message with a one-byte compression flag and a four-byte big-endian length. Specializing \texttt{grpc::SerializationTraits<T>} over \texttt{ctpb::encode} and \texttt{ctpb::decode} is sufficient to carry \texttt{ctpb} types across standard gRPC channels, and automated tests confirm that framed messages are processed correctly by unmodified gRPC server handlers.
